\clearpage\pagestyle{otMain}\setlength{\otBlocEspacement}{2.5mm}\fontsize{11}{15}\selectfont\corrigesChapitre{16}{Triangles et cercle circonscrit}
\begin{blocLivre}{cor-16-1}\begin{corrige}{1}{Le troisième angle}\demarche{Utiliser la somme des angles.}
$\widehat{C}=180-48-67=\res{65^\circ}$.
\label{fig-cor-16-1}\figureGeo[corrige]{tri-1}
\end{corrige}\end{blocLivre}
\begin{blocLivre}{cor-16-2}\begin{corrige}{2}{Triangle isocèle}\demarche{Partager ce qui reste en deux.}
Les angles à la base sont égaux et leur somme vaut $180-40=140^\circ$ : chacun mesure $\res{70^\circ}$.
\label{fig-cor-16-2}\figureGeo[corrige]{tri-2}
\end{corrige}\end{blocLivre}
\begin{blocLivre}{cor-16-3}\begin{corrige}{3}{Triangle équilatéral}\demarche{Trois angles égaux, trois côtés égaux.}
Ses trois angles sont égaux : $180:3=\res{60^\circ}$. Périmètre : $3\times4=\res{12\ \mathrm{cm}}$.
\label{fig-cor-16-3}\figureGeo[corrige]{tri-3}
\end{corrige}\end{blocLivre}
\begin{blocLivre}{cor-16-4}\begin{corrige}{4}{Triangle rectangle}\demarche{Utiliser l’angle droit.}
Les deux angles aigus ont une somme de $90^\circ$ : $\widehat{B}=90-35=\res{55^\circ}$.
\label{fig-cor-16-4}\figureGeo[corrige]{tri-4}
\end{corrige}\end{blocLivre}
\begin{blocLivre}{cor-16-5}\begin{corrige}{5}{Trois longueurs}\demarche{Deux arcs de cercle donnent le troisième sommet.}
On trace [AB], puis l’arc de centre A de rayon 4 cm et l’arc de centre B de rayon 5 cm ; ils se coupent en C. La construction figure en vert.
\label{fig-cor-16-5}\figureGeo[corrige]{tri-5}
\end{corrige}\end{blocLivre}
\begin{blocLivre}{cor-16-6}\begin{corrige}{6}{Possible ou pas ?}\demarche{Comparer le plus grand côté à la somme des deux autres.}
$3+4=7<8$ : les arcs ne se rencontrent pas, c’est $\res{\text{impossible}}$, comme le montre la figure. $5+6=11>10$ : c’est $\res{\text{possible}}$.
\label{fig-cor-16-6}\figureGeo[corrige]{tri-6}
\end{corrige}\end{blocLivre}
\begin{blocLivre}{cor-16-7}\begin{corrige}{7}{Deux côtés et un angle}\demarche{Construire l’angle avant de reporter la longueur.}
On trace [AB], on construit l’angle de $60^\circ$ en A, puis on place C à 3 cm de A sur le second côté. La construction est en vert.
\label{fig-cor-16-7}\figureGeo[corrige]{tri-7}
\end{corrige}\end{blocLivre}
\begin{blocLivre}{cor-16-8}\begin{corrige}{8}{Un côté et deux angles}\demarche{Deux angles et le côté qui les relie.}
Les deux demi-droites tracées depuis B et C se coupent en A (en vert). $\widehat{A}=180-50-60=\res{70^\circ}$.
\label{fig-cor-16-8}\figureGeo[corrige]{tri-8}
\end{corrige}\end{blocLivre}
\begin{blocLivre}{cor-16-9}\begin{corrige}{9}{Angles d’un isocèle}\demarche{Les deux angles à la base sont égaux.}
$\widehat{A}=180-72-72=\res{36^\circ}$.
\label{fig-cor-16-9}\figureGeo[corrige]{tri-9}
\end{corrige}\end{blocLivre}
\begin{blocLivre}{cor-16-10}\begin{corrige}{10}{Construire le cercle circonscrit}\demarche{Deux médiatrices donnent le centre.}
On trace les médiatrices de [AB] et de [BC] ; elles se coupent en O. Le cercle de centre O et de rayon OA passe par A, B et C. La construction est en vert ; la troisième médiatrice passe aussi par O.
\label{fig-cor-16-10}\figureGeo[corrige]{tri-10}
\end{corrige}\end{blocLivre}
\begin{blocLivre}{cor-16-11}\begin{corrige}{11}{Trois villages}\demarche{Traduire chaque équidistance par une médiatrice.}
L’antenne doit être à égale distance de A et de B, et de B et de C : elle est sur deux médiatrices. On la place en O, $\res{\text{centre du cercle circonscrit}}$ au triangle ABC, en vert.
\label{fig-cor-16-11}\figureGeo[corrige]{tri-11}
\end{corrige}\end{blocLivre}
\begin{blocLivre}{cor-16-12}\begin{corrige}{12}{Le cercle d’un triangle rectangle}\demarche{Construire, puis observer.}
Les médiatrices se coupent au $\res{\text{milieu de [BC]}}$ : le centre est sur le côté opposé à l’angle droit, et [BC] est un diamètre du cercle.
\label{fig-cor-16-12}\figureGeo[corrige]{tri-12}
\end{corrige}\end{blocLivre}
\begin{blocLivre}{cor-16-13}\begin{corrige}{13}{Isocèle et équilatéral}\demarche{Enchaîner les propriétés des deux triangles.}
$\widehat{C}=\widehat{B}=65^\circ$, donc $\widehat{BAC}=180-130=\res{50^\circ}$. $\widehat{CAD}=60^\circ$, d’où $\widehat{BAD}=50+60=\res{110^\circ}$.
\label{fig-cor-16-13}\figureGeo[corrige]{tri-13}
\end{corrige}\end{blocLivre}
\begin{blocLivre}{cor-16-14}\begin{corrige}{14}{Dans un triangle équilatéral}\demarche{Utiliser la bissectrice, puis la somme des angles.}
$\widehat{BAH}=60:2=\res{30^\circ}$. Dans le triangle ABH : $\widehat{AHB}=180-60-30=\res{90^\circ}$. La bissectrice est perpendiculaire à [BC].
\label{fig-cor-16-14}\figureGeo[corrige]{tri-14}
\end{corrige}\end{blocLivre}
\begin{blocLivre}{cor-16-15}\begin{corrige}{15}{Deux angles obtus ?}\demarche{Utiliser la somme des angles.}
$100+95=195$, déjà plus que $180^\circ$ : c’est $\res{\text{impossible}}$. La figure le confirme : les deux côtés s’écartent et ne se rencontrent jamais. Un triangle a au plus un angle obtus.
\label{fig-cor-16-15}\figureGeo[corrige]{tri-15}
\end{corrige}\end{blocLivre}
